\chapter{Research lineage, scope, and routes to application}
\label{researchmap:chapter}

\section{The subject and its place in existing science}

Constructive access concerns a concrete experimental question: given a target
property, a class of possible mechanisms, and an apparatus, which conclusions
can be justified, and what additional experiment would justify more?
The apparatus includes preparation, interventions, timing, detection, retained
records, and resource limits. Its mathematical object is the family of
transcript laws produced by the permitted protocols. This formulation supplies
a common language for systems whose underlying physics can be quite different.

This subject lies at the intersection of system identification, statistical
decision theory, controlled stochastic models, quantum characterization, and
thermodynamic inference. Each has a substantial history. Aggregated Markov
models and their indistinguishable realizations were studied in ion-channel
kinetics \cite{p_kienker,p_larget}. Finite linear representations and equivalence
algorithms originate in automata and realization theory
\cite{f_tzeng,f_kiefer,f_thon}. Observable operator models and predictive state
representations already replace inaccessible coordinates by predictions of
future experiments \cite{f_jaeger,f_littman}. A predictive coordinate need not
name a unique molecule, hidden conformation, or physical degree of freedom.

The classical--quantum linear correspondence is also established. Quantum
finite-state equivalence uses finite-dimensional operator spaces
\cite{f_wang,p_faigle}; a recent unified treatment constructs canonical
generalized hidden Markov representations for classical and quantum models
\cite{p_riechers}. The present use of an accessible effect space inherits that
lineage. Similarly, risk-based comparison of experiments and optimal recovery
provide established foundations for deciding which distinctions matter and
how errors limit recovery \cite{f_blackwell,p_donoho}.

The contribution developed in this reference is an explicit chain from those
foundations to a checkable experimental conclusion. A user specifies a
physical access model and a target; obtains an obstruction or a distinguishing
protocol; verifies its algebraic certificate; assigns calibration and sampling
errors; and states the physical resource boundary under which the conclusion
holds. The same chain includes abstention when the supplied model class fails
or the remaining models disagree about the target. The particular model
results and software artifacts are useful parts of that chain. Their status
is recorded separately from the established general principles.

\section{How the earlier results fit together}

The initial boundary question asked when a local observation must display an
arrow of time inherited from driving elsewhere. A finite ring made the
boundary mechanism explicit: one connection carries the driving, while a
binary detector reports a region. This separated three targets that an
informal arrow-of-time argument can conflate: a driven hidden dynamics,
irreversibility of the recorded process, and incompatibility of that recording
with every equilibrium explanation.

In the four-state adjacent-region example, the recording can be reversible
even when the ring is driven; on the specified reciprocal branch an explicit
equilibrium realization produces the same complete recording. Larger-ring
classifications, detector placement, and weak-drive response orders then
showed that access depends on geometry and acquisition. Retaining timestamps
or changing readout spacing can change the leading information order without
changing the hidden physical process. These are exact classifications of the
stated models, rather than statements about the existence of time.

The next step changed access through interventions. In the driven triangle,
each constant-control singleton record is blind to current orientation,
whereas a sequence of two controls produces a current-dependent return
contrast. An autonomous reversible kinetic gate provides a second route: a
binary three-time correlation acquires a leading term proportional to the
controlled edge current and the gate profile's Dirichlet form. The gate's own
steady channel entropy contribution can vanish while the source remains
driven. Preparation, sensing, and control implementation still require their
own accounting.

The quantum branch showed why the target must remain explicit. For the
specified three-site chain, one family of central instruments conceals the
bias magnitude, while a larger family exposes it. Reflection symmetry still
conceals the sign under every permitted central-only protocol. A calibrated
spatial asymmetry changes that conclusion. Magnitude, orientation, and a
complete microscopic model are different identification tasks.

The finite-instrument formulation subsequently provided a common algebraic
interface for these examples. Its cost certificates concern the cheapest
nonzero witness under a prescribed additive charge. Robust testing and
finite-budget decision design concern reliable inference. These objectives
require separate calculations: a very cheap nonzero witness may have much
worse reliable-decision cost than a more informative experiment.

The channel-resolved application studies expose a further distinction between an
inaccessible state and an inaccessible transition channel. Two reaction
mechanisms can have the same full state generator yet different physical
currents and dissipation. Adding the same state-level control preserves that
equivalence. A calibrated channel-selective intervention can remove it in the
worked two-channel model. This is a more precise rule for choosing access
than the suggestion that any temporal modulation reveals hidden dynamics.

Table~\ref{researchmap:lineage} summarizes the retained role of each stage.
The machine-readable provenance record binds individual inherited claims to
their source and result hashes; the table is a reading map, not a count of
discoveries.

\begin{table}[htbp]
\centering\small
\begin{tabular}{p{0.23\textwidth}p{0.67\textwidth}}
\hline
Stage & Retained contribution to the reference\\
\hline
Boundary and detector models & Explicit source, sensor, equilibrium mimics,
sampling effects, and response-order examples.\\[2pt]
Controlled hidden currents & Interventions that alter observable distinction;
control placement and preparation failure cases.\\[2pt]
Autonomous equilibrium gate & Current--correlation identity, triangle
activation criterion, and the cost of hiding the controller record.\\[2pt]
Finite instruments & A common exact representation, finite equivalence
test, witness synthesis, and bounded decision design.\\[2pt]
Certification and adequacy & Independently checked cost certificates,
conditional prediction tests, calibrated abstention, and energy tails.\\[2pt]
Physical diagnostic models & Selective channel identification, targeted
quantum diagnosis, fixed-nuisance sensing, and hidden-energy obstructions.\\
\hline
\end{tabular}
\caption{The dependency structure of the research. Each row supplies a
different part of the experimental argument.}
\label{researchmap:lineage}
\end{table}

The subsequent stages extend this chain in two directions. Compatible-model
results turn evidence into requirements on hidden state count, preparation and
sustained dissipation, with explicit equilibrium countermodels. The virtual
laboratories establish how timing, detector exposure and serial dependence
alter those certificates. The calibrated-inference chapters then improve
statistical use of detector moments, analog records and sequential acquisition,
while proving limits on total-dissipation reconstruction. Their joint
calibration test treats the extra measurement directly as part of the final
scientific contrast. The public-data chapter tests this chain against an
actual observation format and records precisely where calibration remains
insufficient.

\section{What is established, and what is specific to this work}

Three-time correlation tests of nonequilibrium and current inference predate
this work \cite{physics-qian}. Equilibrium kinetic coupling can also carry trajectory
information \cite{p_barato}. The specific gate result here is the factorized
coefficient, its controlled assumptions, the triangle converse, and the stated
information-order comparison. The general-network coefficient supplies a
sufficient activation condition; a vanishing coefficient need not imply
complete invisibility. Historical priority for this precise combination
remains open.

The independent cost checker verifies explicit filtered closure identities,
including those needed to exclude every cheaper word. Its useful distinction
is the exported proof obligation: a user can check the claimed objective
without trusting the search history. Exact equivalence, invariant-space
certificates, and weighted witness methods already exist
\cite{f_kiefer,f_thon}. The implementation does not establish that its supplied
matrices describe an apparatus, nor does it convert arbitrary charges into
joules.

The adequacy layer uses established nonnegative betting processes
\cite{caWaudbySmithRamdas2024,p_intervals}. Its application obligation is to supply valid
conditional prediction envelopes under the actual interventions and histories.
That is stronger than supplying stationary averages. The selective-response
criterion and its deterministic error bounds use linear identifiability and
optimal-recovery geometry. Their exact certificates make the assumptions and
the amplification of calibration error visible.

Recent thermodynamic inference already extracts kinetic information from
partially observed transitions, including periodically driven systems
\cite{p_maier}. Care is needed when applying thermodynamic time reversal after
coarse observation: a physically justified reversal operation and a defined
observation map are part of the model \cite{p_dieball}. This reference's
instantaneous finite-state detectors have an explicit reversal convention;
extensions to milestoned, filtered, or history-dependent detector outputs
must re-establish the relevant correspondence.

Adaptive quantum design also has experimental precedents
\cite{p_stace,physics-bonato}. The worked quantum and Ramsey calculations here
provide transparent small models, explicit alternative hypotheses, and
matched decision objectives. They are reference calculations against which
an implementation can be checked. Demonstrating a hardware advantage
requires comparison with appropriate existing methods on the same device.

\section{A practical standard for extending the subject}

A new application should identify the target before selecting its mathematics.
For a current, the target must name the physical channel and its units. For a
quantum fault, it must name a justified error dictionary and the operations
permitted during diagnosis. For a sensor, it must specify whether the task is
sign discrimination, local parameter precision, tracking, or a decision with
unequal consequences. For energy inference, it must state whether the target
is native-channel free-energy dissipation, gross supplied work, or total
apparatus consumption.

The minimum useful transfer includes a negative control. Examples are a
channel decomposition that produces the same state record, a control family
that preserves a quantum symmetry, an insufficiently selective intervention,
or a missing candidate that forces abstention. A successful case should then
change exactly the access assumption responsible for that failure and predict
the resulting effect before evaluating new data.

Four external milestones would establish practical maturity in a new setting:
\begin{enumerate}
\item An independent specialist checks the theorem, the software contract,
and the closest prior methods for the stated model class.
\item Independent calibration and physical measurements test the preparation,
intervention, detector, memory, and resource assumptions.
\item A comparison holds the scientific objective, error criterion, and full
resource boundary fixed; an increase in contrast alone is insufficient.
\item A second implementation or apparatus reproduces the prediction and
tests at least one intended transfer beyond the development example.
\end{enumerate}
These are measurable acceptance conditions. The present release supplies
mathematical derivations, executable models, and internal adversarial checks.
It also includes the exploratory published-data reanalysis in
Chapter~\ref{ch:ci-empirical}, which identifies concrete apparatus-calibration
gaps. It contains no new laboratory measurements, empirical entropy
validation or external peer-review decisions.

\section{Extensions with a clear scientific payoff}

Several extensions now have a precise starting point. A validated nonlinear
response family can replace the linear selective-response model, provided
identifiability and uncertainty propagation are proved for that family.
Chapter~\ref{ch:continuous} supplies a finite CTMC adapter. Extensions
beyond that finite-model, finite-menu contract to continuous controls or
infinite-dimensional dynamics must also carry approximation error into
the final decision bound, rather than report only solver tolerance. A structured many-body
model can reduce computational burden, provided the structure is physically
justified and the bound includes deviations from it. Detector memory, drift,
and correlated preparation can be handled by enlarging the state or by
conditional envelopes that actually cover those effects.

Chapter~\ref{completion:composition} now supplies a finite classical composition contract and a constructive deterministic interface refinement. Beyond that scoped result, composition remains a concrete research direction. Individually
uninformative components can become informative when the allowed coupling
changes their transcript laws; combining records that remain jointly
equivalent cannot do so. A useful composition theorem must specify the joint
preparation, shared hidden variables, allowed couplings, and target. This
distinguishes a real increase in experimental capability from merely adding
more components or more processing.

The broad claim is therefore an operational one: scientific conclusions can
be organized around what specified experiments preserve, erase, and expose.
The constructive finite theory makes a substantial part of that question
computable and reviewable. Further generality is earned by extending those
interfaces and proving their contracts, with the physical interpretation
checked in each new setting.
